% Standalone manuscript: all text, proofs, and bibliography are in this file.
% Compile twice with pdfLaTeX; no external bibliography or figure files are required.
\documentclass[11pt]{article}
\usepackage[a4paper,margin=1in]{geometry}
\usepackage[T1]{fontenc}
\usepackage[utf8]{inputenc}
\usepackage{lmodern}
\usepackage{microtype}
\usepackage{amsmath,amssymb,amsthm,mathtools}
\usepackage{aliascnt}
\usepackage{booktabs,array,tabularx}
\usepackage{flafter}
\usepackage{enumitem}
\usepackage{xcolor}
\usepackage[authoryear,round]{natbib}
\usepackage{xurl}
\usepackage{needspace}
\definecolor{linkblue}{RGB}{30,66,100}
\usepackage[colorlinks=true,linkcolor=linkblue,citecolor=linkblue,urlcolor=linkblue]{hyperref}
\hypersetup{pdftitle={Smooth Convex Optimization with Simultaneous Euclidean and lp Radius Bounds},pdfauthor={Colombo Team},pdfsubject={First-order oracle complexity of non-Euclidean smooth convex optimization},pdfkeywords={convex optimization, oracle complexity, non-Euclidean smoothness, randomized lower bounds}}
\usepackage[nameinlink,capitalise,noabbrev]{cleveref}
\setlist[itemize]{leftmargin=1.5em,itemsep=3pt,topsep=4pt}
\setlist[enumerate]{leftmargin=1.7em,itemsep=3pt,topsep=4pt}
\setlength{\emergencystretch}{2em}
\setlength{\tabcolsep}{4pt}
\setlength{\parskip}{2pt}
\allowdisplaybreaks[2]
\raggedbottom
\newcolumntype{Y}{>{\raggedright\arraybackslash}X}
\newcommand{\R}{\mathbb{R}}
\newcommand{\E}{\mathbb{E}}
\newcommand{\Prob}{\mathbb{P}}
\newcommand{\N}{\mathbb{N}}
\newcommand{\cF}{\mathcal{F}}
\newcommand{\cA}{\mathcal{A}}
\newcommand{\cO}{\mathcal{O}}
\newcommand{\cS}{\mathcal{S}}
\newcommand{\ip}[2]{\left\langle #1,#2\right\rangle}
\newcommand{\norm}[2]{\left\lVert #1\right\rVert_{#2}}
\newcommand{\abs}[1]{\left|#1\right|}
\newcommand{\argmin}{\operatorname*{arg\,min}}
\newcommand{\vol}{\operatorname{vol}}
\newcommand{\sign}{\operatorname{sign}}
\newcommand{\Risk}{\mathfrak{R}}
\newcommand{\Qdet}{Q_{\mathrm{det}}}
\newcommand{\Qexp}{Q_{\mathrm{E}}}
\newcommand{\Qhp}{Q_{\mathrm{hp}}}
\newcommand{\softO}{\widetilde{O}}
\newcommand{\softOm}{\widetilde{\Omega}}
\newcommand{\softTh}{\widetilde{\Theta}}
\newcommand{\eps}{\varepsilon}
\newcommand{\defeq}{\mathrel{:=}}
\newcommand{\fstar}{f^\star}
\newcommand{\xstar}{x^\star}
\newtheorem{theorem}{Theorem}[section]
\newaliascnt{lemma}{theorem}
\newtheorem{lemma}[lemma]{Lemma}
\aliascntresetthe{lemma}
\newaliascnt{proposition}{theorem}
\newtheorem{proposition}[proposition]{Proposition}
\aliascntresetthe{proposition}
\newaliascnt{corollary}{theorem}
\newtheorem{corollary}[corollary]{Corollary}
\aliascntresetthe{corollary}
\theoremstyle{definition}
\newaliascnt{definition}{theorem}
\newtheorem{definition}[definition]{Definition}
\aliascntresetthe{definition}
\newtheorem{algorithm}{Algorithm}
\theoremstyle{remark}
\newaliascnt{remark}{theorem}
\newtheorem{remark}[remark]{Remark}
\aliascntresetthe{remark}
\crefname{theorem}{Theorem}{Theorems}
\crefname{lemma}{Lemma}{Lemmas}
\crefname{proposition}{Proposition}{Propositions}
\crefname{corollary}{Corollary}{Corollaries}
\crefname{definition}{Definition}{Definitions}
\crefname{remark}{Remark}{Remarks}
\crefname{algorithm}{Algorithm}{Algorithms}
\crefname{section}{Section}{Sections}
\crefname{appendix}{Appendix}{Appendices}
\crefname{equation}{Equation}{Equations}
\title{Smooth Convex Optimization with Simultaneous\\Euclidean and $\ell_p$ Radius Bounds}
\author{Colombo Team\\[3pt]\small\href{mailto:colombo.2026@outlook.com}{\texttt{colombo.2026@outlook.com}}}
\date{}
\begin{document}
\maketitle
\begin{abstract}
We study unconstrained convex minimization when the objective is globally
$L$-smooth in $\ell_p$, $2\le p\le\infty$, and one attained global minimizer
satisfies both a Euclidean radius bound $R_2$ and an $\ell_p$ radius bound
$R_p$. In the unrestricted exact first-order oracle model, the query
complexity is
\[
\widetilde\Theta\!\left(\min\left\{d,\,
\sqrt{\frac{LR_2^2}{\eps}},\,
\left(\frac{LR_p^2}{\eps}\right)^{p/(p+2)}\right\}\right)
\]
throughout the positive-query accuracy range, with $p/(p+2)=1$ at
$p=\infty$. The same rate holds for deterministic error, randomized
expected error, and randomized error with any prescribed failure
probability at most $1/3$. For each fixed finite $p$ and $1\le T\le d/100$,
the expected minimax error is, up to constants depending only on $p$,
$L\min\{R_2^2/(T+1)^2,\,R_p^2/(T+1)^{1+2/p}\}$.
The upper bound combines classical acceleration in two geometries with
finite-dimensional localization. The lower bound uses one distribution
over globally smooth functions, sampled independently of the algorithm.
Its main ingredient controls both the affine pieces made visible by
smoothing and the information conveyed by pieces that remain hidden.
The argument allows unbounded adaptive queries and an arbitrary output,
while every hard instance has an attained minimizer satisfying both
radius bounds. The result identifies how the two classical acceleration
rates combine under a common geometric promise.
\end{abstract}

\noindent\textbf{Keywords.} Convex optimization; first-order oracle complexity;
non-Euclidean smoothness; randomized lower bounds; simultaneous radius bounds.
\par
\medskip
\begin{samepage}
\noindent\textbf{AI Usage.}
Nearly the entire research pipeline for this paper was carried out by
Colombo, powered by GPT-5.6 Sol. 
\par
\end{samepage}


\section{Introduction}
\label{sec:introduction}

Smooth convex optimization has a familiar Euclidean benchmark: with
smoothness constant $L$ and distance $R_2$ to a minimizer, acceleration
achieves error of order $LR_2^2/T^2$ after $T$ first-order queries
\citep{Nesterov1983,NemirovskiiNesterov1985}. For smoothness measured in
$\ell_p$, $p>2$, geometry-adapted methods instead yield the classical
single-radius scale $LR_p^2/T^{1+2/p}$
\citep{GuzmanNemirovski2015,dAspremontGuzmanJaggi2018}.
These formulas involve different distances to a minimizer. Comparing
their exponents alone therefore leaves a basic question unanswered:
what is the optimal rate when both distances are known?

We address this question for globally defined, unconstrained objectives.
After translating the initial point to the origin, assume that
$f:\R^d\to\R$ is convex and
\begin{equation}
\label{eq:intro-smooth}
 \norm{\nabla f(x)-\nabla f(y)}{p^*}
 \le L\norm{x-y}{p}\qquad(x,y\in\R^d),
\end{equation}
where $1/p+1/p^*=1$. We require the existence of a single global minimizer
$x^\star$ satisfying
\begin{equation}
\label{eq:intro-radii}
 \norm{x^\star}{2}\le R_2,
 \qquad \norm{x^\star}{p}\le R_p.
\end{equation}
The simultaneous promise matters because
$\norm{x}{p}\le\norm{x}{2}\le d^{1/2-1/p}\norm{x}{p}$.
An upper bound expressed in $R_2$ and a lower bound expressed in $R_p$
need not refer to the same class of instances.

The oracle returns the exact pair $(f(x),\nabla f(x))$. Algorithms may
query any point in $\R^d$, use arbitrary computation between queries, and
return a point that was never queried. Thus the question concerns the
information available through a first-order oracle, independently of
span restrictions or the cost of implementing auxiliary operations.

\paragraph{Main result.}
Define
\[
 A=\sqrt{\frac{LR_2^2}{\eps}},\qquad
 B=\left(\frac{LR_p^2}{\eps}\right)^{p/(p+2)},\qquad
 S=\min\{d,A,B\},
\]
with the limiting exponent $p/(p+2)=1$ for $p=\infty$.
We prove that
\begin{equation}
\label{eq:intro-frontier}
 Q(\eps)=\widetilde\Theta(S)
 \qquad\left(0<\eps<\frac L2\min\{R_2^2,R_p^2\}\right).
\end{equation}
Here $Q$ may denote deterministic complexity, randomized complexity for
expected error, or randomized complexity with any prescribed failure
probability $\delta\le1/3$. Query budgets are worst-case budgets.
The displayed accuracy range is exact: at and above its upper endpoint,
zero queries suffice; below it, every criterion requires at least one.
The explicit logarithmic factors are given in \cref{thm:main}.

For a fixed finite $p$ and $1\le T\le d/100$, the minimax expected error
satisfies the sharper statement
\begin{equation}
\label{eq:intro-risk}
 \Risk_T^{\mathrm{rand}}
 =\Theta_p\!\left(
 L\min\left\{\frac{R_2^2}{(T+1)^2},
              \frac{R_p^2}{(T+1)^{1+2/p}}\right\}\right).
\end{equation}
The deterministic minimax error has the same order. The subscript $p$
allows constants depending on $p$; these constants are not asserted to
remain bounded as $p\to\infty$.

\paragraph{Contributions and relation to classical rates.}
The exponents in the two acceleration branches are established results.
Our contribution is a characterization of their minimum under the
simultaneous promise \cref{eq:intro-radii}, including the dimension term
and unrestricted randomized algorithms. Three features are central.
First, the hard distribution consists of globally smooth functions with
an attained unconstrained minimum, so the geometric promise is realized
by an actual common minimizer. Second, the randomized lower bound requires
only $d\ge100T$ for the fixed-$p$ risk statement. Third, the information
argument accounts for the positive-width band of affine pieces that a
smooth exact oracle can reveal, as well as the posterior information
left by all unexposed pieces.

The latter point is important for randomized lower bounds. Random
continuous offsets remove exact ties, but a smoothing operation can still
depend on several nearly maximal pieces. We prove a bounded-density
inequality controlling this simultaneous exposure and a companion
inequality controlling the total logarithmic loss of posterior mass.
Together they imply a constant expected number of newly revealed
coordinates per query. This argument works for arbitrary query
magnitudes and charges the final output without requiring another
oracle call.

\paragraph{Proof outline.}
The upper bound selects among three known mechanisms.
Condition \cref{eq:intro-smooth} implies Euclidean $L$-smoothness, giving
the $A$ branch. An accelerated method with a uniformly convex power prox
gives the $B$ branch. A center-of-gravity method on
$B_2(R_2)\cap B_p(R_p)$ gives the dimension-dependent branch; its guarantee
counts oracle calls and does not assert efficient centroid computation.

For the lower bound, choose $m\le d$ and set
$a=\min\{R_2/\sqrt m,R_p/m^{1/p}\}$. Independently sample signs
$\sigma_i\in\{-1,1\}$ and offsets $b_i\in[a/2,a]$, and form
\[
 h(x)=\max\{0,\sigma_1x_1+b_1,\ldots,\sigma_mx_m+b_m\}.
\]
An $\ell_s$ infimal convolution, followed by a smoothness normalization,
produces the hard objective. The point $x_i^\star=-a\sigma_i$ is a zero
of $h$ and remains a minimizer after smoothing. The random
coordinate maximum and local smoothing build on established lower-bound
constructions \citep{GuzmanNemirovski2015,BraunGuzmanPokutta2017}.
The proof here supplies the information accounting needed for this
smooth, simultaneous-radius formulation.

\paragraph{Reading guide.}
\Cref{sec:model-results} specifies the model and states the results.
\Cref{sec:related} compares the guarantees under aligned assumptions.
\Cref{sec:upper,sec:lower} explain the upper and lower bounds, and
\cref{sec:implications} describes the resulting geometric regimes.
The appendices contain the full auxiliary, algorithmic, and
distributional proofs.


\section{Oracle model and main results}
\label{sec:model-results}

All logarithms are natural. We use $1/\infty=0$, $p^*=p/(p-1)$ for finite
$p$, and $\infty^*=1$. We interpret $p/(p+2)=1$ when $p=\infty$.
For $r>0$, $B_p(r)$ denotes the closed $\ell_p$
ball centered at the origin. Universal constants do not depend on
$d,p,L,R_2,R_p,\eps$, or $\delta$.

\begin{definition}[Function class]
\label{def:class}
For $d\ge1$, $2\le p\le\infty$, and $L,R_2,R_p>0$, let
$\cF_p(d,L,R_2,R_p)$ be the class of differentiable convex functions
$f:\R^d\to\R$ such that
\begin{equation}
\label{eq:model-smooth}
 \norm{\nabla f(x)-\nabla f(y)}{p^*}
 \le L\norm{x-y}{p}\qquad\text{for all }x,y\in\R^d,
\end{equation}
and there exists a global minimizer $x^\star$ satisfying
\begin{equation}
\label{eq:model-radii}
 \norm{x^\star}{2}\le R_2,\qquad \norm{x^\star}{p}\le R_p.
\end{equation}
The two inequalities concern the same minimizer. We write
$f^\star=\min_{x\in\R^d}f(x)$.
\end{definition}

The algorithm knows $d,p,L,R_2,R_p$ and accesses $f$ through
$\cO_f(x)=(f(x),\nabla f(x))$. Query points and the final output are
measurable functions of previous responses and, for a randomized
algorithm, an independent seed. They may lie anywhere in $\R^d$.
We place no restriction on intermediate computation. In particular,
the minimizer promise defines a class of functions; it does not
restrict the query domain or require the output to lie in a specified
feasible set.

\begin{definition}[Query-complexity criteria]
\label{def:criteria}
For $\eps>0$, let $\Qdet(\eps)$ be the least integer $T\ge0$ for which a
deterministic algorithm using at most $T$ calls satisfies
\[
 \sup_{f\in\cF_p}[f(\widehat x)-f^\star]\le\eps.
\]
Let $\Qexp(\eps)$ be the corresponding minimum for a randomized algorithm
with
\[
 \sup_{f\in\cF_p}\E[f(\widehat x)-f^\star]\le\eps.
\]
For $0<\delta<1$, let $\Qhp(\eps,\delta)$ be the corresponding minimum
with
\[
 \sup_{f\in\cF_p}\Prob(f(\widehat x)-f^\star>\eps)\le\delta.
\]
The call bound holds for every instance and every realization of the
algorithmic seed. Expectations and probabilities in these definitions
are over that seed alone.
\end{definition}

\begin{proposition}[Exact zero-query threshold]
\label{prop:zero-query}
Set $r=\min\{R_2,R_p\}$. For each of the three criteria, with
$\delta<1/2$ in the probability criterion,
\begin{equation}
\label{eq:zero-query}
 Q=0\quad\Longleftrightarrow\quad\eps\ge\frac{Lr^2}{2}.
\end{equation}
\end{proposition}

\begin{proof}
Because $\nabla f(x^\star)=0$, the smoothness upper model gives
$f(0)-f^\star\le L\norm{x^\star}{p}^2/2\le Lr^2/2$.
For the converse, consider
\begin{equation}
\label{eq:two-quadratics}
 f_\pm(x)=\frac L2(x_1\mp r)^2.
\end{equation}
Both functions belong to $\cF_p$. A zero-query output $X$ is independent
of the choice of sign, and
\[
 \frac12\E[f_+(X)+f_-(X)]
 =\frac L2\E[X_1^2+r^2]\ge\frac{Lr^2}{2}.
\]
This excludes smaller expected error, including deterministic error.
When $\eps<Lr^2/2$, the first-coordinate intervals on which
$f_+\le\eps$ and $f_-\le\eps$ are disjoint. Their success probabilities
sum to at most one, so one instance has failure probability at least
$1/2$.
\end{proof}

Throughout the positive-query regime, retain $r=\min\{R_2,R_p\}$ and define
\begin{equation}
\label{eq:scales}
 H=\log(ed),\qquad
 A=\sqrt{\frac{LR_2^2}{\eps}},\qquad
 B=\left(\frac{LR_p^2}{\eps}\right)^{p/(p+2)},\qquad
 S=\min\{d,A,B\}.
\end{equation}

\begin{theorem}[Complexity under simultaneous radius bounds]
\label{thm:main}
There are universal constants $c,C>0$ such that, for every $d\ge1$,
$2\le p\le\infty$, $L,R_2,R_p>0$, $0<\eps<Lr^2/2$, and
$0<\delta\le1/3$, each of
$\Qdet(\eps)$, $\Qexp(\eps)$, and $\Qhp(\eps,\delta)$ satisfies
\begin{equation}
\label{eq:main-explicit}
 \frac{cS}{H}\le Q\le
 C\left[1+\min\left\{
 A,\ H^2 B,\ d\left(1+\log\frac{Lr^2}{\eps}\right)
 \right\}\right].
\end{equation}
Consequently,
\begin{equation}
\label{eq:main-soft}
 Q=\softTh\!\left(\min\left\{d,\,
 \sqrt{\frac{LR_2^2}{\eps}},\,
 \left(\frac{LR_p^2}{\eps}\right)^{p/(p+2)}\right\}\right).
\end{equation}
The upper bound is achieved deterministically and has no dependence on
$\delta$.
\end{theorem}

The soft notation in \cref{eq:main-soft} hides only logarithmic factors
in $ed$ and $Lr^2/\eps$, as specified in \cref{eq:main-explicit}.
In particular, the dimension branch is
$O(d[1+\log(Lr^2/\eps)])$; it does not assert an accuracy-independent
$O(d)$ bound. For each fixed finite $p$, the factor $H^2$ in the
$B$ branch can be replaced by a constant depending only on $p$.

\begin{definition}[Fixed-horizon minimax error]
\label{def:risks}
For an integer $T\ge0$, define
\begin{align*}
 \Risk_T^{\mathrm{det}}
 &=\inf_{\cA\text{ deterministic}}\sup_{f\in\cF_p}
       [f(\cA^f)-f^\star],\\
 \Risk_T^{\mathrm{rand}}
 &=\inf_{\cA\text{ randomized}}\sup_{f\in\cF_p}
       \E[f(\cA^f)-f^\star],
\end{align*}
where the infima range over algorithms using at most $T$ oracle calls.
\end{definition}

\begin{corollary}[Sharp fixed-$p$ error before dimension saturation]
\label{cor:risk}
For each fixed $2\le p<\infty$, there are $c_p,C_p>0$ such that, for
every integer $1\le T\le d/100$,
\begin{equation}
\label{eq:risk-rate}
\begin{split}
 c_p L\min\left\{\frac{R_2^2}{(T+1)^2},
                 \frac{R_p^2}{(T+1)^{1+2/p}}\right\}
 &\le\Risk_T^{\mathrm{rand}}\le\Risk_T^{\mathrm{det}}\\
 &\le C_p L\min\left\{\frac{R_2^2}{(T+1)^2},
                       \frac{R_p^2}{(T+1)^{1+2/p}}\right\}.
\end{split}
\end{equation}
\end{corollary}

The lower bound follows from the next distributional statement, which
also records the role of the auxiliary smoothing exponent.

\begin{theorem}[A fixed smooth hard distribution]
\label{thm:distribution}
Let $16\le m\le d$ be an integer, let $2\le s\le p$ be finite, and set
\begin{equation}
\label{eq:kappa-def}
 \kappa=(s-1)m^{2/s-2/p}.
\end{equation}
There is a distribution supported on $\cF_p(d,L,R_2,R_p)$ such that
every randomized algorithm making at most $m/100$ exact first-order
queries satisfies
\begin{equation}
\label{eq:distribution-lower}
 \Prob\left(f(\widehat x)-f^\star\ge
 \frac{L}{384\kappa}
 \min\left\{\frac{R_2^2}{m^2},
             \frac{R_p^2}{m^{1+2/p}}\right\}\right)\ge\frac34.
\end{equation}
The probability is over the sampled objective and the independent
algorithmic seed. The distribution is the same for all algorithms;
queries may be unbounded and the output may be any measurable point.
\end{theorem}

Taking $s=p<\infty$ yields the constant-factor lower bound in
\cref{cor:risk}. To treat all $p$ uniformly, take
$s=\min\{p,\max\{2,\log m\}\}$, which gives
\begin{equation}
\label{eq:kappa-log}
 \kappa\le e^2\log(em).
\end{equation}
\Cref{app:main-theorem} performs the inversion and covers small
dimensions using \cref{prop:zero-query}.


\section{Comparison with prior work}
\label{sec:related}

The relevant comparisons involve four separate choices: the smoothness
norm, the geometric promise on the solution, the oracle, and the
dimension required by a lower bound. \Cref{tab:related-comparison}
keeps these choices explicit. The bounds in its first three rows
concern minimization over a norm ball; the last row concerns the
unconstrained class in \cref{def:class}.

\paragraph{Classical acceleration and local smoothing.}
Euclidean acceleration supplies the $R_2$ upper bound
\citep{Nesterov1983,NemirovskiiNesterov1985}; standard treatments of
oracle complexity include \citet{NemirovskiYudin1983,Nesterov2018,Bubeck2015}.
For $p\ge2$, \citet[Corollary~1]{GuzmanNemirovski2015} establish
large-scale deterministic local-oracle lower bounds, using local
smoothing of signed, shifted maxima of affine functions.
Their model already allows oracle queries in the ambient space.
The issue addressed here is the simultaneous-radius promise for an
attained unconstrained minimizer and the corresponding fixed
distribution against randomized algorithms.

Uniformly convex prox functions give the matching single-radius
upper rate; see \citet{dAspremontGuzmanJaggi2018}.
Our power-prox proof is a self-contained specialization of this
established acceleration framework. We use it to make the
dependence on $p$ and the query count explicit, and claim no new
acceleration exponent.

\begin{table}[htbp]
\centering
\caption{First-order function-value guarantees for $p\ge2$.
Here $\alpha=1+2/p$ and $\lambda_{p,d}=\min\{p,\log(ed)\}$.
The symbol $T$ denotes oracle calls, and constants with subscript $p$
may depend on a fixed finite $p$.
The dimension condition in the third row is sufficient for the displayed
constant-confidence specialization; it is not claimed necessary.}
\label{tab:related-comparison}
\small
\renewcommand{\arraystretch}{1.15}
\begin{tabularx}{\textwidth}{@{}>{\raggedright\arraybackslash}p{0.29\textwidth}>{\raggedright\arraybackslash}p{0.35\textwidth}Y@{}}
\toprule
Work and solution geometry & Error guarantee & Algorithms and regime \\
\midrule
\citet{GuzmanNemirovski2015}; minimization over $B_p(R_p)$
& Lower: $\displaystyle\Omega\!\left(\frac{LR_p^2}{\lambda_{p,d}(T+1)^\alpha}\right)$
& Deterministic local-oracle algorithms; $1\le T$, $T+1\le d$; $2\le p\le\infty$. \\
\addlinespace[5pt]
\citet{dAspremontGuzmanJaggi2018}; minimization over $B_p(R_p)$
& Upper: $\displaystyle O_p\!\left(\frac{LR_p^2}{T^\alpha}\right)$
& Geometry-adapted acceleration; fixed $2\le p<\infty$. The exponent is classical. \\
\addlinespace[5pt]
\citet[Theorem~14]{DiakonikolasGuzman2020}; one norm-ball geometry
& Lower: $\displaystyle\Omega_p\!\left(\frac{LR_p^2}{T^\alpha}\right)$
& Randomized local-oracle methods; fixed finite $p$, constant confidence; sufficient $d\gtrsim_p T^{3+2/p}\log(eT)$. \\
\addlinespace[5pt]
This paper, \cref{cor:risk}; one unconstrained minimizer obeys both bounds
& Matching: $\displaystyle\Theta_p\!\left(L\min\left\{\frac{R_2^2}{(T+1)^2},\frac{R_p^2}{(T+1)^\alpha}\right\}\right)$
& Deterministic and randomized expected error; exact value and gradient; fixed finite $p$, $1\le T\le d/100$. \\
\bottomrule
\end{tabularx}
\end{table}

For the Guzm\'an--Nemirovski bound, an arbitrary output is appended as
one further query, accounting for $T+1$ in the displayed rate and
dimension condition.

\paragraph{Randomized and parallel lower bounds.}
\citet{DiakonikolasGuzman2020} provide a general framework for
randomized and parallel local-oracle lower bounds in normed spaces.
The third row of \cref{tab:related-comparison} follows from their
smooth case with one query per round: setting
$\eps\asymp_p LR_p^2T^{-(1+2/p)}$ in their dimension condition gives
the sufficient requirement displayed there.
\citet{ContrerasGuzmanMartinezRubio2025} extend non-Euclidean
optimization to higher-order smoothness and include randomized and
parallel lower bounds; their first-order specialization recovers the
same single-radius exponent. Their results cover more oracle orders
and parallelism than the sequential first-order model studied here.

For nonsmooth objectives, \citet{BraunGuzmanPokutta2017}, together
with their published correction \citep{BraunGuzmanPokutta2024Correction},
develop information arguments based on string guessing, signed
coordinate maxima, and continuous perturbations. These are direct
antecedents of our construction. Our technical task is to control the
entire band visible after smoothing and the cumulative information in
the inequalities satisfied by unexposed pieces.

\paragraph{Trajectory-sensitive acceleration.}
For the same standard $\ell_p$ smoothness condition,
\citet[Theorem~4]{BaiBullins2025} obtain
\begin{equation}
\label{eq:intro-hasd}
 f(x_T)-f^\star\le
 \frac{324L\norm{x_0-x^\star}{2}^2}{\mathcal G_T^2T^2},
 \qquad
 \mathcal G_T=\frac1T\sum_{t=0}^{T-1}
 \frac{\norm{\nabla f(x_{t+1})}{p^*}}
      {\norm{\nabla f(x_{t+1})}{2}}.
\end{equation}
Here $T$ counts outer iterations; their binary-search implementation
uses logarithmically many oracle calls per iteration
\citep[Theorem~5]{BaiBullins2025}. A zero gradient permits termination
at a minimizer. The quantity $\mathcal G_T$ depends on the realized
trajectory. It can explain substantial gains on individual instances,
but deriving a uniform class improvement from this guarantee would
require a corresponding lower bound on $\mathcal G_T$ over the
entire class.

\paragraph{Standard and adaptive smoothness.}
\citet{XieWangWuLi2026} distinguish standard norm smoothness from
adaptive smoothness based on a common matrix majorant of the Hessian.
For diagonal preconditioners, their parameters satisfy
$L_\infty\le\Lambda_{\mathrm{diag}}\le dL_\infty$
\citep[Proposition~2.5]{XieWangWuLi2026}.
An acceleration bound controlled by $\Lambda_{\mathrm{diag}}$ therefore
has a different normalization from one controlled by $L_\infty$.
The present class uses only \cref{eq:model-smooth} and the two distances
of one global minimizer.


\section{Three deterministic upper bounds}
\label{sec:upper}

The parameters of the class are known, so an algorithm can select the
method with the smallest certified call bound before making its first
query. This section records the three guarantees in a common model.
The acceleration proofs and the full localization argument appear in
\cref{app:upper-proofs}.

\subsection{Euclidean acceleration}

Since $p^*\le2\le p$, norm monotonicity gives
\begin{equation}
\label{eq:p-implies-euclidean}
 \norm{\nabla f(x)-\nabla f(y)}{2}
 \le\norm{\nabla f(x)-\nabla f(y)}{p^*}
 \le L\norm{x-y}{p}
 \le L\norm{x-y}{2}.
\end{equation}
Every function in $\cF_p$ is thus Euclidean $L$-smooth with the same
constant. Classical acceleration yields the following guarantee.

\begin{proposition}[Euclidean branch]
\label{prop:euclidean}
For every integer $T\ge1$, a deterministic method using one gradient
query per iteration returns $x_T$ with
\begin{equation}
\label{eq:euclidean-error}
 f(x_T)-f^\star\le\frac{2LR_2^2}{(T+1)^2}.
\end{equation}
In particular, $C(1+A)$ calls suffice for error at most $\eps$.
\end{proposition}

The same class inclusion proves an upper bound, but transferring a
lower bound requires separate care: a hard Euclidean-smooth function
need not satisfy \cref{eq:model-smooth} with the same $L$.

\subsection{Acceleration with a power prox}

For any finite $2\le q\le p$, condition \cref{eq:model-smooth} also
implies $L$-smoothness in $\ell_q$, because
$\norm{g}{q^*}\le\norm{g}{p^*}$ and $\norm{x}{p}\le\norm{x}{q}$.
The minimizer obeys
\begin{equation}
\label{eq:Rq}
 \norm{x^\star}{q}\le d^{1/q-1/p}R_p\defeq R_q.
\end{equation}
We use the classical power-prox construction
\citep{dAspremontGuzmanJaggi2018}, in the explicit normalization
\begin{equation}
\label{eq:power-prox}
 \psi_q(z)=\frac{2^{q-2}}{q}\norm{z}{q}^q.
\end{equation}
Its Bregman divergence is at least $\norm{u-v}{q}^q/q$.
Combining this uniform convexity with an inexact order-$q$ model gives
a direct estimate-sequence proof.

\begin{proposition}[Power-prox branch]
\label{prop:power}
Suppose $f$ is convex and $L$-smooth in $\ell_q$, where $2\le q<\infty$,
and has a global minimizer satisfying $\norm{x^\star}{q}\le R$.
A deterministic method using one gradient query per iteration satisfies
\begin{equation}
\label{eq:power-query}
 T\ge4q^2\left(\frac{LR^2}{\eps}\right)^{q/(q+2)}
 \quad\Longrightarrow\quad f(x_T)-f^\star\le\eps.
\end{equation}
For $q=2$, the method also satisfies
$f(x_T)-f^\star\le2LR^2/(T+1)^2$.
\end{proposition}

For fixed finite $p$, take $q=p$ to obtain $C_p(1+B)$ calls.
For a bound uniform in $p$, including $p=\infty$, take
\begin{equation}
\label{eq:q-cap}
 q=\min\{p,\max\{2,\log d\}\}.
\end{equation}
Then
\begin{equation}
\label{eq:q-cap-properties}
 d^{2/q-2/p}\le e^2,\qquad q\le2\log(ed)=2H.
\end{equation}
Indeed, either $q=p$, or $q\ge\log d$.
In the positive-query range $LR_p^2/\eps>2$.
Using \cref{eq:Rq} and the monotonicity of $u/(u+2)$, a sufficient
integer horizon in \cref{eq:power-query} is bounded by
\begin{equation}
\label{eq:power-uniform}
 C\left[1+H^2
 \left(\frac{LR_p^2}{\eps}\right)^{p/(p+2)}\right]
 =C(1+H^2B).
\end{equation}

\subsection{Finite-dimensional localization}

Both promises identify a known convex body containing $x^\star$:
\begin{equation}
\label{eq:known-body}
 K=B_2(R_2)\cap B_p(R_p).
\end{equation}
At the centroid $c$ of its current localization body, the algorithm
queries $f$ and retains the halfspace
\begin{equation}
\label{eq:gradient-cut}
 \ip{\nabla f(c)}{x-c}\le0.
\end{equation}
Convexity keeps every global minimizer inside the body.
Gr\"unbaum's inequality \citep{Grunbaum1960} ensures a constant-factor
volume reduction whenever the gradient is nonzero.

\begin{proposition}[Dimension branch]
\label{prop:dimension}
For $0<\eps<Lr^2/2$, there is a deterministic algorithm using at most
\begin{equation}
\label{eq:dimension-query}
 Cd\left(1+\log\frac{Lr^2}{\eps}\right)
\end{equation}
oracle calls and returning a queried point with error at most $\eps$.
\end{proposition}

The key volume comparison remains valid even if the chosen minimizer
lies on the boundary of $K$. Set
\[
 \alpha=\sqrt{\frac{\eps}{2Lr^2}},\qquad
 K_\eps=(1-\alpha)x^\star+\alpha K.
\]
Every $x\in K_\eps$ satisfies $\norm{x-x^\star}{p}\le2\alpha r$,
so the smoothness upper model gives $f(x)-f^\star\le\eps$.
If every queried centroid has error greater than $\eps$, all of
$K_\eps$ survives every cut. Its volume is
$\alpha^d\vol(K)$, contradicting repeated constant-factor volume
reduction after the number of calls in \cref{eq:dimension-query}.
The algorithm returns the queried point of smallest objective value;
it need not know $f^\star$.

Exact centroid computation is permitted by the oracle model and may
be expensive. \Cref{prop:dimension} therefore bounds information
requirements, while \cref{prop:euclidean,prop:power} provide explicit
first-order iterations. Selecting the smallest of their certified
budgets proves the upper side of \cref{thm:main} for all three criteria.


\section{A fixed distribution for the randomized lower bound}
\label{sec:lower}

We now describe the construction underlying \cref{thm:distribution}.
All random parameters are sampled before the algorithm runs. The proof
compares the exact first-order oracle with an enhanced oracle that
reveals additional information. A lower bound for this more informative
interaction then applies to the original oracle.

\subsection{Local smoothing that preserves a minimum}

For finite $s\ge2$ and a nonnegative convex function $h$ that is
$1$-Lipschitz in $\ell_s$, define
\begin{equation}
\label{eq:smoothing}
 (\cS_\mu h)(x)
 =\min_y\left\{h(y)+\frac{\norm{x-y}{s}^2}{2\mu}\right\}.
\end{equation}
This norm-adapted infimal convolution is a standard local smoothing
mechanism; see \citet{GuzmanNemirovski2015}. We record the precise
properties needed here, including preservation of a zero minimum.

\begin{lemma}[Local smoothing]
\label{lem:smoothing}
Let $h:\R^m\to[0,\infty)$ be convex and $1$-Lipschitz in $\ell_s$,
where $2\le s<\infty$. Then
\begin{align}
 0\le\cS_\mu h\le h,\qquad
 \cS_\mu h&\ge h-\mu/2,
 \label{eq:smoothing-approx}\\
 \norm{y_h(x)-x}{s}&\le\mu,
 \label{eq:smoothing-locality}
\end{align}
where $y_h(x)$ is the unique minimizer in \cref{eq:smoothing}.
The envelope is differentiable and $(s-1)/\mu$-smooth in $\ell_s$.
If $h(x^\star)=0$, then $(\cS_\mu h)(x^\star)=0$.
If two functions satisfying these assumptions agree on
$B_s(x,2\mu)$, their envelope values and gradients agree at $x$.
\end{lemma}

\Cref{app:smoothing} gives the full proof. The displacement bound makes
the oracle local, and nonnegativity preserves an attained unconstrained
minimum without an additional penalty or a feasible-set boundary.

\subsection{The sampled hard instance}

Fix the parameters of \cref{thm:distribution} and set
\begin{equation}
\label{eq:a-mu-tau}
 a=\min\left\{\frac{R_2}{\sqrt m},\frac{R_p}{m^{1/p}}\right\},
 \qquad\mu=\frac{a}{128m},\qquad\tau=4\mu=\frac{a}{32m}.
\end{equation}
Independently sample
\begin{equation}
\label{eq:latent}
 \sigma_i\sim\operatorname{Unif}\{-1,1\},\qquad
 b_i\sim\operatorname{Unif}[a/2,a],\qquad i=1,\ldots,m.
\end{equation}
For $u\in\R^m$, let
\begin{equation}
\label{eq:hard-max}
 h(u)=\max\{0,\sigma_1u_1+b_1,\ldots,\sigma_mu_m+b_m\},
\end{equation}
and for $x\in\R^d$ define
\begin{equation}
\label{eq:hard-smooth}
 f(x)=\frac{L\mu}{\kappa}(\cS_\mu h)(x_1,\ldots,x_m).
\end{equation}
Signed coordinate maxima with random perturbations have been used in
information-based lower bounds \citep{BraunGuzmanPokutta2017}.
Here the nonnegative maximum and the scale $a$ enforce the attained
minimum and both radius constraints.

Specifically, choose
\begin{equation}
\label{eq:hard-minimizer}
 x_i^\star=-a\sigma_i\quad(i\le m),\qquad
 x_i^\star=0\quad(i>m).
\end{equation}
All coordinate-affine pieces are nonpositive there, so
$h(x_1^\star,\ldots,x_m^\star)=0$ and
$f(x^\star)=0=\min f$. Moreover,
\begin{equation}
\label{eq:hard-radii}
 \norm{x^\star}{2}=a\sqrt m\le R_2,\qquad
 \norm{x^\star}{p}=am^{1/p}\le R_p.
\end{equation}
Each affine piece has $\ell_{s^*}$ slope norm at most one, hence $h$ is
$1$-Lipschitz in $\ell_s$. Since
$\norm{v}{s}\le m^{1/s-1/p}\norm{v}{p}$, the envelope's quadratic upper
model has coefficient
$(s-1)m^{2/s-2/p}/\mu=\kappa/\mu$ in $\ell_p$.
The scaling in \cref{eq:hard-smooth} therefore makes $f$ globally
$L$-smooth in $\ell_p$; \cref{lem:norm-conversion} supplies the formal
gradient-norm implication. Every sampled objective belongs to the
stated class.

\subsection{What a smooth oracle can reveal}

At a query $x$, write
\begin{equation}
\label{eq:Y-and-M}
 Y_i=\sigma_i x_i+b_i,\qquad
 M=\max\{0,Y_1,\ldots,Y_m\}.
\end{equation}
Knowing a single maximizer need not determine a smoothed gradient.
The enhanced oracle therefore reveals every previously hidden pair
$(\sigma_i,b_i)$ in the band
\begin{equation}
\label{eq:near-band}
 Y_i\ge M-\tau.
\end{equation}
For all other hidden coordinates, it reports the inequality
\begin{equation}
\label{eq:posterior-truncation}
 Y_i<M-\tau.
\end{equation}
Each such inequality truncates a coordinate's posterior support.
Whenever that support has prior mass less than $3/4$, the oracle also
reveals the pair. We call these two mechanisms \emph{band exposure}
and \emph{forced exposure}, respectively.

Conditional on this transcript and the algorithmic seed, the remaining
pairs are independent restrictions of their original priors, each to
a set of prior mass at least $3/4$. Thus either sign has conditional
probability at most $2/3$. Also, for every possible query coordinate
$x_i$, the conditional density of $Y_i$ obeys
\begin{equation}
\label{eq:density-bound-main}
 D_i\le\frac{8}{3a}.
\end{equation}
Before normalization, each sign branch contributes density at most
$1/a$; the support restriction can increase the total density by at
most $4/3$. This bound is independent of the magnitude of $x_i$.
\Cref{lem:posterior} proves the conditional product statement.

The enhanced response determines the exact value and gradient of
$f$. Every omitted piece is more than $4\mu$ below an included
maximizer at $x$. On $B_s(x,2\mu)$, the difference between two
coordinate-affine pieces changes by at most $4\mu$. Omitted pieces
are therefore inactive throughout this ball. Locality in
\cref{lem:smoothing} gives the exact oracle simulation
(\cref{lem:simulation}).

\subsection{Exposure and posterior information}

The following lemma controls both kinds of information acquired from
one query.

\begin{lemma}[Near-maximum exposure and posterior loss]
\label{lem:band}
Let $Y_1,\ldots,Y_n$ be independent real random variables whose
continuous distribution functions $F_i$ have densities bounded by
$D_i$. For a fixed $c\in\R$ and $\tau>0$, let
\[
 M=\max\{c,Y_1,\ldots,Y_n\},\qquad
 J=\{i:Y_i\ge M-\tau\}.
\]
Then
\begin{align}
 \E\abs{J}&\le1+\tau\sum_iD_i,
 \label{eq:band-count}\\
 \E\sum_{i\notin J}-\log F_i(M-\tau)
 &\le1+\tau\sum_iD_i.
 \label{eq:band-loss}
\end{align}
The same bounds hold conditionally on any history preserving the
stated independence and density bounds. For $n=0$, both left-hand
sides are zero.
\end{lemma}

The first inequality counts the possible maximizer and then bounds
the chance that another variable falls in an interval of width $\tau$.
The second inequality is the additional accounting needed here.
Writing $F=\prod_iF_i$ and $Z=\max_iY_i$, the total loss at the
unshifted maximum is at most $-\log F(M)$, whose expectation is at
most one because $F(Z)$ is uniform. Moving the cutoff down by $\tau$
costs at most $\tau\sum_iD_i$: condition on the maximum of the other
variables and use
$v\log(w/v)\le w-v$. The complete proof is in
\cref{app:revelation}.

For the hard distribution,
\begin{equation}
\label{eq:tauD}
 \tau\sum_iD_i\le\frac{a}{32m}\frac{8m}{3a}=\frac1{12}.
\end{equation}
Consequently, one query has at most $13/12$ expected band exposures
and incurs at most $13/12$ expected units of logarithmic posterior
mass loss. A forced exposure consumes more than $\log(4/3)$ units of
accumulated loss from its original prior. Charging that loss only
once per exposed coordinate yields the following bound.

\begin{lemma}[Expected exposure budget]
\label{lem:budget}
Let $R_T$ be the number of exposed coordinates after $T$ adaptive
enhanced-oracle calls. Then
\begin{equation}
\label{eq:budget}
 \E R_T\le
 \frac{13}{12}\left(1+\frac1{\log(4/3)}\right)T<5T.
\end{equation}
The expectation includes the independent algorithmic seed.
\end{lemma}

\subsection{Converting hidden signs to error}

If $T\le m/100$, Markov's inequality gives
\begin{equation}
\label{eq:half-hidden}
 \Prob(R_T>m/2)\le\frac1{10}.
\end{equation}
Otherwise at least $m/2$ signs remain independent conditional on the
transcript, and neither value of any remaining sign has probability
greater than $2/3$. Once the transcript and seed are fixed, the output
$\widehat x$ is fixed as well. To have
$\sigma_i\widehat x_i<0$, each nonzero output coordinate must select
the correct sign; a zero coordinate cannot satisfy the strict
inequality. Thus, for $m\ge16$,
\begin{equation}
\label{eq:sign-error}
 \Prob(\sigma_i\widehat x_i<0\text{ for every }i\le m)
 \le\frac1{10}+\left(\frac23\right)^{m/2}<\frac14.
\end{equation}
With probability at least $3/4$, some coordinate therefore gives
$h(\widehat x_1,\ldots,\widehat x_m)\ge b_i\ge a/2$.
By \cref{eq:smoothing-approx}, its smoothed value is at least $a/3$.
It follows that
\[
 f(\widehat x)-f^\star
 \ge\frac{L\mu a}{3\kappa}
 =\frac{L}{384\kappa}
 \min\left\{\frac{R_2^2}{m^2},
             \frac{R_p^2}{m^{1+2/p}}\right\}.
\]
This is \cref{thm:distribution}. \Cref{app:revelation} supplies the
complete posterior, simulation, and exposure proofs.

The order of quantifiers is worth making explicit: one distribution
is fixed first, and the conclusion then holds for every randomized
algorithm. Averaging over that distribution gives a fixed hard
instance for each algorithm. The proof neither selects a sign after
observing the seed nor infers a randomized lower bound from a
deterministic resisting oracle.


\section{Geometric regimes and consequences}
\label{sec:implications}

\subsection{Redundant radii and the crossover}

Norm comparison identifies the range in which both radius parameters
carry distinct information. More precisely, replacing the radii by
\begin{equation}
\label{eq:canonical-radii}
 R_2'=\min\{R_2,d^{1/2-1/p}R_p\},\qquad
 R_p'=\min\{R_p,R_2\}
\end{equation}
leaves $B_2(R_2)\cap B_p(R_p)$, and hence the function class, unchanged.
Indeed, every point in the original intersection already satisfies
both tightened bounds. The reverse inclusion follows because the
new radii are no larger. One may therefore assume
\begin{equation}
\label{eq:nonredundant-radii}
 R_p\le R_2\le d^{1/2-1/p}R_p.
\end{equation}
The theorem itself applies without this normalization.

For $2<p<\infty$, equating the two fixed-horizon rate functions in
\cref{eq:risk-rate} gives the crossover scale
\begin{equation}
\label{eq:general-crossover}
 T_{\mathrm{cross}}\asymp
 \left(\frac{R_2}{R_p}\right)^{2p/(p-2)},\qquad
 \eps_{\mathrm{cross}}\asymp
 LR_2^2\left(\frac{R_p}{R_2}\right)^{4p/(p-2)}.
\end{equation}
Here replacing $T+1$ by $T$ affects only constants. For horizons below
this scale, the smaller $R_p$ can compensate for the slower exponent
$1+2/p$. At longer horizons the Euclidean $T^{-2}$ branch becomes
smaller, provided the horizon still lies in the range covered by
\cref{cor:risk}. Under \cref{eq:nonredundant-radii}, the formal crossover
lies between order one and order $d$.

At $p=2$, both exponents are two and the smaller radius governs the
rate. At $p=\infty$, the analogous polynomial scales are
$T_{\mathrm{cross}}\asymp(R_2/R_\infty)^2$ and
$\eps_{\mathrm{cross}}\asymp LR_\infty^4/R_2^2$;
the theorem at this endpoint retains logarithmic factors.

\subsection{An explicit example: \texorpdfstring{$p=4$}{p=4}}

Set $L=R_2=1$ and write $R_4=\rho$. The nonredundant radius range is
$d^{-1/4}\le\rho\le1$. The query bound becomes
\begin{equation}
\label{eq:p4-query}
 Q(\eps)=\softTh\!\left(
 \min\{d,\eps^{-1/2},\rho^{4/3}\eps^{-2/3}\}\right),
 \qquad 0<\eps<\rho^2/2.
\end{equation}
For $1\le T\le d/100$,
\begin{equation}
\label{eq:p4-risk}
 \Risk_T^{\mathrm{rand}}\asymp\Risk_T^{\mathrm{det}}
 \asymp\min\{(T+1)^{-2},\rho^2(T+1)^{-3/2}\}.
\end{equation}
Thus
\begin{equation}
\label{eq:p4-crossover}
 T_{\mathrm{cross}}\asymp\rho^{-4},\qquad
 \eps_{\mathrm{cross}}\asymp\rho^8.
\end{equation}

\begin{table}[htbp]
\centering
\caption{Polynomial query regimes for $p=4$ and $L=R_2=1$.
The rows compare the three terms in \cref{eq:p4-query}; logarithmic
factors and constant-width transitions are suppressed.
Always restrict to $0<\eps<\rho^2/2$; an empty interval contributes
no regime.}
\label{tab:p4}
\small
\renewcommand{\arraystretch}{1.3}
\begin{tabularx}{\textwidth}{@{}p{0.28\textwidth}p{0.29\textwidth}Y@{}}
\toprule
Accuracy scale & Smallest polynomial term & Mechanism \\
\midrule
$\rho^8\lesssim\eps<\rho^2/2$
& $\rho^{4/3}\eps^{-2/3}$ & Acceleration in $\ell_4$ exploits the smaller radius. \\
\addlinespace[3pt]
$d^{-2}\lesssim\eps\lesssim\rho^8$
& $\eps^{-1/2}$ & Euclidean acceleration benefits from the faster decay in $T$. \\
\addlinespace[3pt]
$0<\eps\lesssim d^{-2}$
& $d$ & Finite-dimensional localization, with an accuracy logarithm. \\
\bottomrule
\end{tabularx}
\end{table}

For example, $\rho=d^{-1/8}$ gives
$T_{\mathrm{cross}}\asymp d^{1/2}$ and
$\eps_{\mathrm{cross}}\asymp d^{-1}$.
For sufficiently large $d$, both acceleration branches are visible before the polynomial
dimension threshold $\eps\asymp d^{-2}$.
At $\rho=1$, the Euclidean branch is smaller throughout the nontrivial
accuracy range. At $\rho=d^{-1/4}$, the crossover reaches the dimension
scale, so the Euclidean intermediate regime collapses.

\subsection{Trajectory gains and worst-case guarantees}

With $R_2=R_4=1$, \cref{eq:p4-risk} gives
\begin{equation}
\label{eq:no-uniform-gain}
 \Risk_T^{\mathrm{rand}}=\Theta(T^{-2}),\qquad 1\le T\le d/100.
\end{equation}
Consequently, standard $\ell_4$ smoothness and these two radii alone
cannot guarantee an additional factor $d^{1/2}$ in every accelerated
error bound. This conclusion concerns the entire oracle model.
The trajectory statistic in \cref{eq:intro-hasd} can still be large on
particular instances; the minimax class also contains instances
preventing a uniform improvement.

\subsection{Randomization and confidence}

The hard distribution gives a positive error threshold with probability
at least $3/4$. The same event yields an expected-error lower bound
and rules out failure probability at most $1/3$ below that threshold.
Conversely, the upper bounds are deterministic. Thus increasing the
required success probability within the stated range introduces no
confidence-amplification cost in the upper bound. This observation
does not identify the optimal logarithmic factors in dimension or
accuracy.


\section{Discussion}
\label{sec:discussion}

The simultaneous-radius class has a simple complexity structure:
two classical acceleration mechanisms compete, and finite-dimensional
localization supplies a third option. The lower bound shows that their
minimum is unavoidable up to logarithmic factors even for randomized
algorithms with unrestricted exact first-order access. For fixed
finite $p$ and $T\le d/100$, the error rate matches up to constants.

The main technical contribution is the information argument for a
locally smoothed random maximum. Direct band exposures and the
posterior constraints on all other coordinates must be accounted for
together. The bounded-density lemma makes this accounting independent
of the magnitudes of adaptive queries. Preserving a zero minimum
during smoothing then places the entire hard distribution inside the
unconstrained class with a common minimizer satisfying both radii.

Several limits of the result remain relevant. The uniform theorem
does not determine the best logarithmic dependence on $d$ or
$\eps$, especially as $p\to\infty$. The localization branch permits
arbitrary internal computation and gives no polynomial-time guarantee
for exact centroid evaluation. The call budgets are worst-case
budgets; an expected-stopping-time criterion would require a separate
statement and argument. Finally, noisy or finite-precision oracles,
strong convexity, and gradient-norm objectives define different
complexity questions.

These distinctions also suggest the next questions: which logarithmic
losses are intrinsic, which computationally efficient localization
methods preserve the dimension dependence, and which additional
geometric promises convert favorable trajectory behavior into a
uniform improvement over the simultaneous-radius minimax rate?


\clearpage
% Standalone natbib-compatible bibliography. Optional recent tangential citations omitted.
\begin{thebibliography}{99}

\bibitem[Bai and Bullins(2025)]{BaiBullins2025}
Cedar Site Bai and Brian Bullins.
Faster acceleration for steepest descent.
In \emph{Proceedings of the Thirty-Eighth Conference on Learning Theory},
volume~291 of \emph{Proceedings of Machine Learning Research}, pages 202--230, 2025.
\url{https://proceedings.mlr.press/v291/bai25a.html}.

\bibitem[Braun et~al.(2017)Braun, Guzm\'an, and Pokutta]{BraunGuzmanPokutta2017}
G\'abor Braun, Crist\'obal Guzm\'an, and Sebastian Pokutta.
Lower bounds on the oracle complexity of nonsmooth convex optimization via information theory.
\emph{IEEE Transactions on Information Theory}, 63(7):4709--4724, 2017.
\url{https://doi.org/10.1109/TIT.2017.2701343}.

\bibitem[Braun et~al.(2024)Braun, Guzm\'an, and Pokutta]{BraunGuzmanPokutta2024Correction}
G\'abor Braun, Crist\'obal Guzm\'an, and Sebastian Pokutta.
Corrections to ``Lower bounds on the oracle complexity of nonsmooth convex optimization via information theory.''
\emph{IEEE Transactions on Information Theory}, 70(7):5408--5409, 2024.
\url{https://doi.org/10.1109/TIT.2024.3357200}.

\bibitem[Bubeck(2015)]{Bubeck2015}
S\'ebastien Bubeck.
Convex optimization: Algorithms and complexity.
\emph{Foundations and Trends in Machine Learning}, 8(3--4):231--357, 2015.
\url{https://doi.org/10.1561/2200000050}.

\bibitem[Contreras et~al.(2025)Contreras, Guzm\'an, and Mart\'inez-Rubio]{ContrerasGuzmanMartinezRubio2025}
Juan Pablo Contreras, Crist\'obal Guzm\'an, and David Mart\'inez-Rubio.
Non-Euclidean high-order smooth convex optimization: Extended abstract.
In \emph{Proceedings of the Thirty-Eighth Conference on Learning Theory},
volume~291 of \emph{Proceedings of Machine Learning Research}, page 1330, 2025.
Full version: \url{https://arxiv.org/abs/2411.08987}.

\bibitem[d'Aspremont et~al.(2018)d'Aspremont, Guzm\'an, and Jaggi]{dAspremontGuzmanJaggi2018}
Alexandre d'Aspremont, Crist\'obal Guzm\'an, and Martin Jaggi.
Optimal affine-invariant smooth minimization algorithms.
\emph{SIAM Journal on Optimization}, 28(3):2384--2405, 2018.
\url{https://doi.org/10.1137/17M1116842}.

\bibitem[Diakonikolas and Guzm\'an(2020)]{DiakonikolasGuzman2020}
Jelena Diakonikolas and Crist\'obal Guzm\'an.
Lower bounds for parallel and randomized convex optimization.
\emph{Journal of Machine Learning Research}, 21(5):1--31, 2020.
\url{https://jmlr.org/papers/v21/19-771.html}.

\bibitem[Gr\"unbaum(1960)]{Grunbaum1960}
Branko Gr\"unbaum.
Partitions of mass-distributions and of convex bodies by hyperplanes.
\emph{Pacific Journal of Mathematics}, 10(4):1257--1261, 1960.
\url{https://doi.org/10.2140/pjm.1960.10.1257}.

\bibitem[Guzm\'an and Nemirovski(2015)]{GuzmanNemirovski2015}
Crist\'obal Guzm\'an and Arkadi Nemirovski.
On lower complexity bounds for large-scale smooth convex optimization.
\emph{Journal of Complexity}, 31(1):1--14, 2015.
\url{https://doi.org/10.1016/j.jco.2014.08.003}.

\bibitem[Nemirovski and Yudin(1983)]{NemirovskiYudin1983}
Arkadi S. Nemirovski and David B. Yudin.
\emph{Problem Complexity and Method Efficiency in Optimization}.
Wiley-Interscience, New York, 1983.

\bibitem[Nemirovskii and Nesterov(1985)]{NemirovskiiNesterov1985}
Arkadi S. Nemirovskii and Yurii E. Nesterov.
Optimal methods of smooth convex minimization.
\emph{USSR Computational Mathematics and Mathematical Physics}, 25(2):21--30, 1985.
\url{https://doi.org/10.1016/0041-5553(85)90100-4}.

\bibitem[Nesterov(1983)]{Nesterov1983}
Yurii E. Nesterov.
A method of solving a convex programming problem with convergence rate $O(1/k^2)$.
\emph{Soviet Mathematics Doklady}, 27:372--376, 1983.

\bibitem[Nesterov(2018)]{Nesterov2018}
Yurii E. Nesterov.
\emph{Lectures on Convex Optimization}.
Springer, 2018.
\url{https://doi.org/10.1007/978-3-319-91578-4}.

\bibitem[Xie et~al.(2026)Xie, Wang, Wu, and Li]{XieWangWuLi2026}
Shuo Xie, Tianhao Wang, Beining Wu, and Zhiyuan Li.
A tale of two geometries: Adaptive optimizers and non-Euclidean descent.
In \emph{International Conference on Learning Representations}, 2026.
\href{https://proceedings.iclr.cc/paper_files/paper/2026/hash/7c79d0964c8fc3674dcf1737e8afe89f-Abstract-Conference.html}{ICLR 2026 proceedings}.

\end{thebibliography}

\clearpage
\appendix
\crefalias{section}{appendix}

\section{Norm inequalities and accelerated weights}
\label{app:preliminaries}

We collect the elementary estimates used in the upper and lower bounds.
All logarithms are natural.

\begin{lemma}[Finite-dimensional norm comparison]
\label{lem:norm-comparison}
For $1\le a\le b\le\infty$ and $x\in\R^m$,
\begin{equation}
\label{eq:norm-comparison}
 \norm{x}{b}\le\norm{x}{a}
 \le m^{1/a-1/b}\norm{x}{b}.
\end{equation}
\end{lemma}

\begin{proof}
The first inequality is norm monotonicity.  For $b<\infty$, the second
follows by applying H\"older's inequality to $\sum_j|x_j|^a$.
The case $b=\infty$ follows directly from
$\sum_j|x_j|^a\le m\norm{x}{\infty}^a$.
\end{proof}

\begin{lemma}[A quadratic upper model implies gradient Lipschitzness]
\label{lem:upper-model-implies-smooth}
Let $F$ be differentiable and convex on a finite-dimensional normed space,
and let $\beta>0$.  If
\begin{equation}
\label{eq:abstract-upper-model}
 F(x+h)\le F(x)+\ip{\nabla F(x)}{h}
             +\frac\beta2\norm{h}{}^2
 \qquad\text{for all }x,h,
\end{equation}
then
\begin{equation}
\label{eq:gradient-Lipschitz-from-model}
 \norm{\nabla F(x)-\nabla F(y)}{*}
 \le\beta\norm{x-y}{}
 \qquad\text{for all }x,y.
\end{equation}
\end{lemma}

\begin{proof}
Fix $x,y$ and set $G(z)=F(z)-\ip{\nabla F(y)}{z}$.
Convexity and $\nabla G(y)=0$ imply that $y$ minimizes $G$.
Applying \cref{eq:abstract-upper-model} to $G$ at $x$, then minimizing the
right-hand side over $h$, gives
\[
 G(y)\le G(x)-\frac{1}{2\beta}\norm{\nabla G(x)}{*}^2.
\]
Here we used the dual-norm identity
$\inf_h\{\ip{g}{h}+\beta\norm{h}{}^2/2\}
=-\norm{g}{*}^2/(2\beta)$.
Thus
\begin{equation}
\label{eq:coercive-bregman}
 F(x)-F(y)-\ip{\nabla F(y)}{x-y}
 \ge\frac{1}{2\beta}
       \norm{\nabla F(x)-\nabla F(y)}{*}^2.
\end{equation}
Adding the same inequality with $x$ and $y$ interchanged yields
\[
 \ip{\nabla F(x)-\nabla F(y)}{x-y}
 \ge\frac1\beta\norm{\nabla F(x)-\nabla F(y)}{*}^2.
\]
H\"older's inequality proves \cref{eq:gradient-Lipschitz-from-model}.
\end{proof}

\begin{lemma}[Smoothness of the squared $\ell_s$ norm]
\label{lem:squared-ls-smooth}
For $2\le s<\infty$, the function
\begin{equation}
\label{eq:phi-s}
 \phi_s(z)=\frac12\norm{z}{s}^2
\end{equation}
satisfies
\begin{equation}
\label{eq:phi-upper-model}
 \phi_s(z+v)\le\phi_s(z)+\ip{\nabla\phi_s(z)}{v}
                    +\frac{s-1}{2}\norm{v}{s}^2.
\end{equation}
Its gradient is therefore $(s-1)$-Lipschitz from $\ell_s$ to $\ell_{s^*}$.
\end{lemma}

\begin{proof}
For $s=2$, the upper model is an identity.  For $s>2$ and $z\ne0$,
direct differentiation gives
\begin{align}
 D^2\phi_s(z)[v,v]
 &=(s-1)\norm{z}{s}^{2-s}
       \sum_j|z_j|^{s-2}v_j^2
 \notag\\
 &\quad-(s-2)\norm{z}{s}^{2-2s}
       \left(\sum_j|z_j|^{s-2}z_jv_j\right)^2
 \notag\\
 &\le(s-1)\norm{z}{s}^{2-s}
       \sum_j|z_j|^{s-2}v_j^2
 \le(s-1)\norm{v}{s}^2.
 \label{eq:phi-hessian}
\end{align}
The last inequality is H\"older's inequality.  Integration along a line
segment proves \cref{eq:phi-upper-model}.  If the segment meets the
origin, apply the same argument on its two sides and use continuity of
$\phi_s$ and its gradient at zero.  The gradient bound follows from
\cref{lem:upper-model-implies-smooth}.
\end{proof}

\begin{lemma}[Accelerated weight growth]
\label{lem:weight-growth}
Let $q\ge2$, $M>0$, and $A_0=0$.  For $t\ge1$, define $a_t>0$ and $A_t$ by
\begin{equation}
\label{eq:weight-equation-appendix}
 A_t=A_{t-1}+a_t,
 \qquad a_t^q=\frac{A_t^{q-1}}{M}.
\end{equation}
The positive solution exists and is unique.  For every integer $T\ge1$,
\begin{equation}
\label{eq:weight-growth}
 A_T\ge\frac{(T+q-1)^q}{q^qM}
       \ge\frac{T^q}{q^qM}.
\end{equation}
In particular, when $q=2$ and $M=L$,
$A_T\ge(T+1)^2/(4L)$.
\end{lemma}

\begin{proof}
For fixed $A_{t-1}\ge0$, the function
$a\mapsto a^q/(A_{t-1}+a)^{q-1}$ is continuous and strictly increasing
from zero to infinity on $(0,\infty)$.  This proves existence and
uniqueness.  The first weight is $A_1=a_1=1/M$.
Concavity of $u\mapsto u^{1/q}$ gives, for $t\ge2$,
\[
 A_t^{1/q}-A_{t-1}^{1/q}
 \ge\frac{A_t-A_{t-1}}{qA_t^{(q-1)/q}}
 =\frac{1}{qM^{1/q}}.
\]
Summing from $t=2$ to $T$ and using $A_1^{1/q}=M^{-1/q}$ proves
\cref{eq:weight-growth}.
\end{proof}

\section{Complete proofs of the upper bounds}
\label{app:upper-proofs}

We first give one algorithm that proves both acceleration bounds.
We then establish the dimension-dependent localization bound.

\subsection{Power-prox acceleration}

\begin{lemma}[Power uniform convexity]
\label{lem:power-uniform-convexity}
For $q\ge2$, the function
\[
 \psi_q(z)=\frac{2^{q-2}}q\norm{z}{q}^q
\]
satisfies
\begin{equation}
\label{eq:power-uc}
 \psi_q(v)\ge\psi_q(u)+\ip{\nabla\psi_q(u)}{v-u}
                       +\frac1q\norm{v-u}{q}^q.
\end{equation}
\end{lemma}

\begin{proof}
For real $a,b$,
\begin{equation}
\label{eq:scalar-monotonicity}
 \bigl(|a|^{q-2}a-|b|^{q-2}b\bigr)(a-b)
 \ge2^{2-q}|a-b|^q.
\end{equation}
Indeed, write the difference on the left as the integral of
$(q-1)|r|^{q-2}$ over the interval joining $a$ and $b$.
Among intervals of a fixed length, this integral is smallest on the
interval centered at zero, which gives the displayed constant.
Multiplying by $2^{q-2}$ and summing over coordinates yields
\[
 \ip{\nabla\psi_q(v)-\nabla\psi_q(u)}{v-u}
 \ge\norm{v-u}{q}^q.
\]
Apply this inequality to $u$ and $u+t(v-u)$, divide by $t>0$, and
integrate over $0<t<1$.  The resulting integral of $t^{q-1}$ is $1/q$,
which proves \cref{eq:power-uc}.
\end{proof}

\begin{lemma}[Inexact power model]
\label{lem:inexact-power-model}
Suppose $f$ is convex and $L$-smooth in $\ell_q$, where $q>2$.
For every $\zeta>0$,
\begin{equation}
\label{eq:inexact-power-model}
 f(x+h)\le f(x)+\ip{\nabla f(x)}{h}
                  +\frac Mq\norm{h}{q}^q+\frac\zeta2,
\end{equation}
where
\begin{equation}
\label{eq:M-power-model}
 M=L^{q/2}\left(\frac{q-2}{q\zeta}\right)^{(q-2)/2}.
\end{equation}
For $q=2$, the same model holds with $M=L$ and no additive error.
\end{lemma}

\begin{proof}
The smoothness remainder is at most $Lt^2/2$, where $t=\norm{h}{q}$.
For $q>2$,
\[
 \max_{t\ge0}\left\{\frac L2t^2-\frac Mqt^q\right\}
 =\frac{q-2}{2q}L^{q/(q-2)}M^{-2/(q-2)}.
\]
Substitution of \cref{eq:M-power-model} makes this maximum $\zeta/2$.
The case $q=2$ is the usual smoothness upper model.
\end{proof}

\par\medskip\noindent
\begin{minipage}{\textwidth}
\begin{algorithm}[Power-prox acceleration with horizon $T\ge1$]
\label{alg:power-prox}
\mbox{}\par\smallskip
\noindent\textbf{Input:} finite $q\ge2$, smoothness $L$, accuracy $\eps>0$, and horizon $T$.
\begin{enumerate}[leftmargin=1.7em,itemsep=4pt]
\item If $q>2$, set $\zeta=\eps/T$ and choose $M$ by \cref{eq:M-power-model}; otherwise set $M=L$.
\item Initialize $A_0=0$, $y_0=z_0=0$, and $G_0=0$.
\item For $t=1,\ldots,T$, perform the following updates:
\begin{enumerate}[label=(\alph*),leftmargin=1.7em,itemsep=4pt]
 \item Choose $a_t>0$ such that $A_t=A_{t-1}+a_t$ and $a_t^q=A_t^{q-1}/M$.
 \item Set $x_t=(A_{t-1}y_{t-1}+a_tz_{t-1})/A_t$.
 \item Query $g_t=\nabla f(x_t)$ and set $G_t=G_{t-1}+a_tg_t$.
 \item Set $(z_t)_j=-\sign((G_t)_j)(|(G_t)_j|/2^{q-2})^{1/(q-1)}$ for every $j$.
 \item Set $y_t=x_t+(a_t/A_t)(z_t-z_{t-1})$.
\end{enumerate}
\item Return $y_T$.
\end{enumerate}
\end{algorithm}
\end{minipage}
\par\medskip

We use the convention $\sign(0)=0$.
The algorithm makes exactly one first-order call per iteration.
Its analysis uses the potential
\begin{equation}
\label{eq:power-potential}
 \Psi_t(z)=\psi_q(z)+
 \sum_{i=1}^t a_i\bigl(f(x_i)+\ip{g_i}{z-x_i}\bigr).
\end{equation}
The coordinate update in \cref{alg:power-prox} is precisely
$z_t=\argmin_z\Psi_t(z)$.

\begin{lemma}[Estimate-sequence invariant]
\label{lem:power-invariant}
For $q>2$, define
\begin{equation}
\label{eq:power-error-accumulation}
 E_t=\frac\zeta2\sum_{i=1}^t A_i.
\end{equation}
Then
\begin{equation}
\label{eq:power-invariant}
 \Psi_t(z_t)\ge A_t f(y_t)-E_t.
\end{equation}
For $q=2$, the same assertion holds with $E_t=0$.
\end{lemma}

\begin{proof}
The assertion holds at $t=0$.  Suppose it holds at $t-1$ and write
$\Delta z_t=z_t-z_{t-1}$.
Linear terms do not change a Bregman divergence.  Thus
\cref{lem:power-uniform-convexity} and the optimality of $z_{t-1}$ give
\begin{align*}
 \Psi_t(z_t)
 &\ge\Psi_{t-1}(z_{t-1})+\frac1q\norm{\Delta z_t}{q}^q
       +a_t\bigl(f(x_t)+\ip{g_t}{z_t-x_t}\bigr)\\
 &\ge A_{t-1}f(y_{t-1})-E_{t-1}
       +\frac1q\norm{\Delta z_t}{q}^q
       +a_t\bigl(f(x_t)+\ip{g_t}{z_t-x_t}\bigr).
\end{align*}
By convexity,
$f(y_{t-1})\ge f(x_t)+\ip{g_t}{y_{t-1}-x_t}$.
Using $A_t x_t=A_{t-1}y_{t-1}+a_tz_{t-1}$ therefore gives
\begin{equation}
\label{eq:power-combine-linearizations}
 \Psi_t(z_t)\ge A_t f(x_t)+a_t\ip{g_t}{\Delta z_t}
                    +\frac1q\norm{\Delta z_t}{q}^q-E_{t-1}.
\end{equation}
Apply \cref{lem:inexact-power-model} with
$h=(a_t/A_t)\Delta z_t$.  The weight equation implies
\[
 A_t\frac Mq\left(\frac{a_t}{A_t}\right)^q=\frac1q.
\]
Since $y_t=x_t+h$, the right-hand side of
\cref{eq:power-combine-linearizations} is at least
\[
 A_t f(y_t)-E_{t-1}-\frac{A_t\zeta}{2}
 =A_t f(y_t)-E_t.
\]
For $q=2$, the same calculation has no additive model error.
\end{proof}

\begin{proof}[Proof of \cref{prop:power,prop:euclidean}]
Each affine model in \cref{eq:power-potential} is a global lower bound
on $f$.  Consequently,
\begin{equation}
\label{eq:power-compare-star}
 \Psi_T(z_T)\le\Psi_T(x^\star)
              \le\psi_q(x^\star)+A_Tf^\star.
\end{equation}
Combining this inequality with \cref{lem:power-invariant} and
$\sum_{t=1}^T A_t\le T A_T$ gives
\begin{equation}
\label{eq:power-gap-general}
 f(y_T)-f^\star\le\frac{\psi_q(x^\star)}{A_T}
                         +\frac{T\zeta}{2}
 \qquad(q>2).
\end{equation}
For $q>2$, use $\psi_q(x^\star)\le2^{q-2}R^q/q$,
\cref{lem:weight-growth}, and $\zeta=\eps/T$ to obtain
\begin{equation}
\label{eq:power-gap-expanded}
 f(y_T)-f^\star
 \le\frac{2^{q-2}q^{q-1}MR^q}{T^q}+\frac\eps2.
\end{equation}
After substituting \cref{eq:M-power-model}, the first term is at most
$\eps/2$ whenever
\[
 T\ge K_q\left(\frac{LR^2}{\eps}\right)^{q/(q+2)},
 \qquad
 K_q=\left[
 2^{q-1}q^{q-1}
 \left(\frac{q-2}{q}\right)^{(q-2)/2}
 \right]^{2/(q+2)}.
\]
Since $(q-2)/q\le1$ and $(q-1)/(q+2)\le1$,
\[
 K_q\le 2^{2(q-1)/(q+2)}q^{2(q-1)/(q+2)}
       \le4q^2.
\]
This proves the claimed power-prox query bound for $q>2$.

For $q=2$, there is no accumulated error.  The sharper weight estimate
$A_T\ge(T+1)^2/(4L)$ and $\psi_2(x^\star)\le R^2/2$ give
\[
 f(y_T)-f^\star\le\frac{2LR^2}{(T+1)^2}.
\]
This proves both the $q=2$ case of \cref{prop:power} and
\cref{prop:euclidean}, using the Euclidean smoothness implication
\cref{eq:p-implies-euclidean} and $R=R_2$ for the latter.
\end{proof}

\subsection{Finite-dimensional localization}

\begin{proof}[Proof of \cref{prop:dimension}]
Work in the positive-query range $0<\eps<Lr^2/2$, where
$r=\min\{R_2,R_p\}$.
Let $K_0=K=B_2(R_2)\cap B_p(R_p)$.
At step $t$, query the centroid $c_t$ of $K_{t-1}$.
If $\nabla f(c_t)=0$, return $c_t$; otherwise set
\[
 K_t=K_{t-1}\cap
 \{x:\ip{\nabla f(c_t)}{x-c_t}\le0\}.
\]
Convexity gives
\[
 \ip{\nabla f(c_t)}{x^\star-c_t}
 \le f(x^\star)-f(c_t)\le0,
\]
so every minimizer in $K$ is retained.
Each nontrivial centroid cut leaves a full-dimensional convex body.
Gr\"unbaum's inequality~\citep{Grunbaum1960} yields
\[
 \vol(K_t)\le(1-e^{-1})\vol(K_{t-1}).
\]

Suppose the first $N$ queried centroids all have error greater than
$\eps$.  Set
\[
 \alpha=\sqrt{\frac{\eps}{2Lr^2}},
 \qquad K_\eps=(1-\alpha)x^\star+\alpha K.
\]
Our accuracy range gives $0<\alpha<1$, and convexity gives
$K_\eps\subseteq K$.
For $x=(1-\alpha)x^\star+\alpha y$ with $y\in K$,
\[
 \norm{x-x^\star}{p}
 =\alpha\norm{y-x^\star}{p}
 \le2\alpha\min\{R_2,R_p\}=2\alpha r.
\]
Indeed, the two available bounds are
$\norm{y-x^\star}{p}\le2R_p$ and
$\norm{y-x^\star}{p}\le\norm{y-x^\star}{2}\le2R_2$.
Because $x^\star$ is an unconstrained minimizer,
$\nabla f(x^\star)=0$, and the smoothness upper model gives
\[
 f(x)-f^\star\le\frac L2\norm{x-x^\star}{p}^2\le\eps.
\]
Thus every $x\in K_\eps$ survives every cut:
\[
 \ip{\nabla f(c_t)}{x-c_t}
 \le f(x)-f(c_t)<0.
\]
It follows that
\[
 \alpha^d\vol(K)=\vol(K_\eps)
 \le\vol(K_N)
 \le(1-e^{-1})^N\vol(K).
\]
This is impossible if
\[
 N>\frac{d\log(1/\alpha)}{-\log(1-e^{-1})}
   =\frac{d\log(2Lr^2/\eps)}{-2\log(1-e^{-1})}.
\]
An integer $N\le C d(1+\log(Lr^2/\eps))$ therefore suffices, for a
universal constant $C$.  Returning the queried centroid of smallest
objective value proves the proposition.  The remaining accuracy range
is handled without queries by \cref{prop:zero-query}.
\end{proof}

\section{Local smoothing and conversion of norms}
\label{app:smoothing}

We prove the smoothing properties needed to construct globally defined
hard instances while preserving an attained minimum.

\begin{proof}[Proof of \cref{lem:smoothing}]
Fix $x\in\R^m$.  The function
\[
 y\longmapsto h(y)+\frac{\norm{x-y}{s}^2}{2\mu}
\]
is coercive and strictly convex, so it has a unique minimizer $y_h(x)$.
Testing $y=x$ gives $\cS_\mu h(x)\le h(x)$, and both terms in the
objective are nonnegative.  Since $h$ is $1$-Lipschitz in $\ell_s$,
\[
 h(y)\ge h(x)-\norm{x-y}{s}.
\]
Therefore
\[
 (\cS_\mu h)(x)
 \ge h(x)+\inf_{t\ge0}\left\{-t+\frac{t^2}{2\mu}\right\}
 =h(x)-\frac\mu2,
\]
proving \cref{eq:smoothing-approx}.

Let $\phi_s(z)=\norm{z}{s}^2/2$.
Optimality of $y_h(x)$ gives
\begin{equation}
\label{eq:prox-optimality-app}
 \frac1\mu\nabla\phi_s(x-y_h(x))\in\partial h(y_h(x)).
\end{equation}
Every subgradient of a convex $1$-Lipschitz function has dual norm at
most one, and direct differentiation yields
\[
 \norm{\nabla\phi_s(z)}{s^*}=\norm{z}{s}.
\]
Thus $\norm{x-y_h(x)}{s}\le\mu$, which proves
\cref{eq:smoothing-locality}.

The infimal convolution is convex.  Its minimizers are locally bounded
in $x$ by the preceding locality estimate.  Uniqueness and Danskin's
theorem therefore give differentiability and the formula
\begin{equation}
\label{eq:envelope-gradient-app}
 \nabla(\cS_\mu h)(x)
 =\frac1\mu\nabla\phi_s(x-y_h(x)).
\end{equation}
Holding $y_h(x)$ fixed and applying \cref{lem:squared-ls-smooth}, we obtain
\begin{align*}
 (\cS_\mu h)(x+v)
 &\le h(y_h(x))+\frac1\mu\phi_s(x+v-y_h(x))\\
 &\le(\cS_\mu h)(x)
       +\ip{\nabla(\cS_\mu h)(x)}{v}
       +\frac{s-1}{2\mu}\norm{v}{s}^2.
\end{align*}
By \cref{lem:upper-model-implies-smooth}, the envelope is
$(s-1)/\mu$-smooth in $\ell_s$.
If $h(x^\star)=0$, the already established inequalities give
$0\le(\cS_\mu h)(x^\star)\le h(x^\star)=0$, so the minimum is preserved.

Finally, suppose $h$ and $\widetilde h$ agree on $B_s(x,2\mu)$.
Their proximal minimizers at $x$ both belong to $B_s(x,\mu)$.
Evaluating each proximal objective at the other's minimizer shows that
the two minimum values are equal.  Strict convexity then implies that
the proximal minimizers coincide.  The gradient formula
\cref{eq:envelope-gradient-app} proves equality of the envelope gradients
as well as their values at $x$.
The same reasoning applies at every point sufficiently close to $x$,
since its proximal minimizers remain in the ball on which the two
functions agree.
\end{proof}

\begin{lemma}[Changing the smoothing norm]
\label{lem:norm-conversion}
Let $2\le s\le p\le\infty$, with $s<\infty$.
If a differentiable convex function $F:\R^m\to\R$ satisfies
\[
 F(x+h)\le F(x)+\ip{\nabla F(x)}{h}
                   +\frac\beta2\norm{h}{s}^2,
\]
then $F$ is $\beta m^{2/s-2/p}$-smooth in $\ell_p$.
\end{lemma}

\begin{proof}
By \cref{lem:norm-comparison},
$\norm{h}{s}\le m^{1/s-1/p}\norm{h}{p}$.
Substitute this bound into the upper model and apply
\cref{lem:upper-model-implies-smooth} in the $\ell_p$ norm.
\end{proof}

For the envelope $\cS_\mu h$, the resulting smoothness coefficient is
\[
 \frac{s-1}{\mu}m^{2/s-2/p}=\frac\kappa\mu.
\]
Multiplication by $L\mu/\kappa$ therefore gives an $L$-smooth function
in $\ell_p$.  Extending this function to $\R^d$ by ignoring the last
$d-m$ coordinates does not increase the smoothness constant: the
$\ell_p$ norm of the first $m$ coordinates is no larger than the full
$\ell_p$ norm, and the remaining gradient coordinates vanish.

\section{Revelation bounds and the fixed hard distribution}
\label{app:revelation}

When evaluating $h$ or $\cS_\mu h$ at an ambient vector in this appendix,
we use its first $m$ coordinates. This appendix gives the complete information argument behind
\cref{thm:distribution}.  Throughout, the signs and offsets are sampled
once from \cref{eq:latent}, independently of the algorithm.  We include
all algorithmic randomness, including any randomness used to produce the
final output, in one independent seed.  Conditioning on this seed and
the past transcript makes each query and the final output measurable
functions of the observed history.  All conditional-distribution
statements below are understood almost surely; exact observations of
continuous offsets are interpreted using regular conditional
distributions.

\subsection{The enhanced oracle and its conditional law}

For each active coordinate, let
\[
 \Omega_i=\{-1,1\}\times[a/2,a],
\]
and let $\nu_i$ denote the prior law of $(\sigma_i,b_i)$ from
\cref{eq:latent}.  The oracle records the exact pair for each exposed
coordinate.  For every unexposed coordinate, it records a measurable
set $E_i\subseteq\Omega_i$ and its original prior mass
\[
 P_i=\nu_i(E_i).
\]
The invariant is that, conditional on the transcript, the unexposed
pairs are independent with respective laws $\nu_i(\cdot\mid E_i)$ and
$P_i\ge3/4$.  Initially, $E_i=\Omega_i$ and $P_i=1$.

At a query $x$, define the known baseline
\[
 c=\max\bigl(\{0\}\cup
       \{\sigma_jx_j+b_j:j\text{ is already exposed}\}\bigr).
\]
For each unexposed index, set $Y_i=\sigma_ix_i+b_i$, and let
$M$ be the maximum of $c$ and these hidden scores.  If there are no
unexposed indices, set $M=c$ and return without changing the state.
Otherwise, the response has the following two steps.
\begin{enumerate}
 \item Reveal the exact pair for every unexposed index satisfying
 $Y_i\ge M-\tau$.  These are the new band exposures.
 \item For every remaining index, update
 \begin{equation}
 \label{eq:E-update}
 E_i\leftarrow
 E_i\cap\{(\sigma,b):\sigma x_i+b<M-\tau\}.
 \end{equation}
 If the updated mass $P_i$ is below $3/4$, reveal its exact pair as a
 forced exposure.  Otherwise, keep the coordinate unexposed with this
 updated set.
\end{enumerate}
The response also reports the updated sets and masses.  These quantities
are determined by the previous state, the query, and the maximum $M$;
granting them does not weaken the lower bound.  The maximum itself is
known after the first step: it is either the baseline $c$ or the score
of a newly exposed pair.

\begin{lemma}[Posterior product structure and density control]
\label{lem:posterior}
After every response, the unexposed latent pairs are conditionally
independent with laws $\nu_i(\cdot\mid E_i)$.  Before each query, either
sign of an unexposed pair has conditional probability at most $2/3$.
Moreover, the conditional density of
$Y_i=\sigma_ix_i+b_i$ is at most $8/(3a)$, uniformly over the query point.
\end{lemma}

\begin{proof}
The initial law is a product.  Suppose that the invariant holds before
a query, and condition on the seed and the past transcript.  The query
and the old sets $E_i$ are then fixed, and the hidden pairs are
independent.

Condition further on the indices and exact pairs exposed in the band.
Their values, together with the baseline, determine $M$.  For every
index outside the band, the only additional restriction on its latent
pair is the coordinatewise event in \cref{eq:E-update}.  In particular,
its failure to appear in the band imposes no further joint restriction
once $M$ is known.  The conditional density of the remaining pairs
therefore factors into their separately restricted densities.

The set of forced exposures is determined by the updated masses, which
are already fixed by the old sets, the query, and $M$.  Revealing the
pairs at those indices preserves independence of all pairs that remain
unexposed.  Their new laws are precisely
$\nu_i(\cdot\mid E_i)$, and their masses are at least $3/4$ by the
exposure rule.  This proves the invariant by induction.

Each sign branch has original prior mass $1/2$.  Consequently,
\[
 \Prob(\sigma_i=1\mid E_i)\le\frac{1/2}{3/4}=\frac23,
 \qquad
 \Prob(\sigma_i=-1\mid E_i)\le\frac23.
\]
The joint prior density of a specified sign and its offset is $1/a$.
For a fixed sign, the map $b\mapsto\sigma x_i+b$ is a translation, so
it contributes density at most $1/a$ before conditioning.  Adding the
two sign branches and dividing by $P_i\ge3/4$ gives
\[
 D_i\le\frac{2/a}{3/4}=\frac{8}{3a}.
\]
The bound does not depend on $x_i$; in particular, it is unaffected by
unbounded query magnitudes.
\end{proof}

\subsection{Simulation of the exact smooth oracle}

\begin{lemma}[Enhanced-oracle simulation]
\label{lem:simulation}
At each query, the enhanced response determines the exact value and
gradient of $\cS_\mu h$, and therefore of the function $f$ in
\cref{eq:hard-smooth}.
\end{lemma}

\begin{proof}
Let $\widetilde h$ be the maximum of zero and all affine pieces exposed
by the end of the current response.  At the query point $x$, one of
these included pieces attains $M$.  Every omitted piece
$\ell_i(z)=\sigma_i z_i+b_i$ satisfies
\[
 \ell_i(x)<M-\tau.
\]
Choose an included piece $\ell_j$ attaining $M$, with
$\ell_j\equiv0$ if the zero piece is selected.  When both pieces are
coordinate-affine, every $z\in B_s(x,2\mu)$ satisfies
\begin{align*}
 \ell_i(z)-\ell_j(z)
 &\le\ell_i(x)-\ell_j(x)
       +\abs{z_i-x_i}+\abs{z_j-x_j}\\
 &< -\tau+4\mu=0.
\end{align*}
If $\ell_j\equiv0$, the possible increase is at most $2\mu$, so the
same strict inequality holds.  Thus no omitted piece is active in
$B_s(x,2\mu)$, and
\[
 h=\widetilde h\qquad\text{on }B_s(x,2\mu).
\]
Both functions are convex, nonnegative, and $1$-Lipschitz in $\ell_s$.
The locality statement in \cref{lem:smoothing} gives equality of their
envelope values and gradients at $x$.  Since $\widetilde h$ is fully
determined by exposed pairs, this identifies the exact oracle response.
The ambient coordinates beyond $m$ are ignored by $f$ and have zero
gradient.
\end{proof}

\subsection{A bound on band exposures and posterior loss}

\begin{proof}[Proof of \cref{lem:band}]
If there are no variables, both sums in the lemma are zero.  Assume
henceforth that $n\ge1$, and define
\[
 M_{-i}=\max\bigl(\{c\}\cup\{Y_j:j\ne i\}\bigr).
\]
Because the variables have continuous distributions and are independent,
there is almost surely at most one unknown variable attaining the
maximum, and ties with $c$ have probability zero.  Every other index
$i\in J$ satisfies
\[
 M_{-i}-\tau\le Y_i\le M_{-i}.
\]
Conditional on $M_{-i}$, the probability of this event is at most
$\tau D_i$, since $Y_i$ and $M_{-i}$ are independent.  Adding one for
the possible unknown maximizer and summing proves
\cref{eq:band-count}.

For the loss bound, let
\[
 Z=\max_iY_i,\qquad F(t)=\prod_iF_i(t).
\]
By independence, $F$ is the continuous CDF of $Z$.  For every
$i\notin J$, decompose
\[
 -\log F_i(M-\tau)
 =-\log F_i(M)
   +\log\frac{F_i(M)}{F_i(M-\tau)}.
\]
These quantities are well defined almost surely on the event under
consideration.  Zero-probability cases may be assigned their limiting
values.  The sum of the first terms is bounded by
\[
 \sum_{i\notin J}-\log F_i(M)
 \le-\log\prod_iF_i(M)
 =-\log F(M)
 \le-\log F(Z),
\]
where the last inequality uses $M\ge Z$.  The probability integral
transform for a continuous CDF gives $F(Z)\sim\operatorname{Unif}[0,1]$;
strict monotonicity of $F$ is not required.  Therefore,
\[
 \E[-\log F(M)]\le\int_0^1-\log u\,du=1.
\]

For the remaining terms, observe that
\[
 \{i\notin J\}=\{Y_i<M_{-i}-\tau\},
\]
and $M=M_{-i}$ on this event.  Conditional on $M_{-i}$, write
$u=M_{-i}-\tau$.  Independence shows that the conditional expected
contribution of index $i$ is
\[
 F_i(u)\log\frac{F_i(u+\tau)}{F_i(u)}.
\]
For $0<v\le w$, the inequality
$v\log(w/v)\le w-v$ implies
\[
 F_i(u)\log\frac{F_i(u+\tau)}{F_i(u)}
 \le F_i(u+\tau)-F_i(u)
 \le\tau D_i.
\]
When $F_i(u)=0$, the expected contribution is zero, and the displayed
bound holds by continuity.  Summing over $i$ and adding the bound of one
for the first terms proves \cref{eq:band-loss}.  Applying this same
argument to conditional laws proves the conditional version of the
lemma.
\end{proof}

\subsection{The cumulative revelation budget}

\begin{proof}[Proof of \cref{lem:budget}]
Condition on the seed and the history before a query.  By
\cref{lem:posterior}, each hidden score has density at most $8/(3a)$.
There are at most $m$ hidden scores, so
\[
 \tau\sum_iD_i
 \le\frac{a}{32m}\frac{8m}{3a}
 =\frac1{12}.
\]
The conditional expected number of new band exposures is therefore at
most $13/12$ by \cref{eq:band-count}.  This remains true when all
coordinates have already been exposed, in which case the count is zero.

For an index outside the band, let $P_i^{\mathrm{old}}$ and
$P_i^{\mathrm{new}}$ be its original prior masses before and after the
truncation.  The conditional CDF before the query satisfies
\[
 F_i(M-\tau)=\frac{P_i^{\mathrm{new}}}{P_i^{\mathrm{old}}}.
\]
The strict truncation and the non-strict CDF agree because the
conditional score distribution is continuous.  The corresponding
log-mass expenditure is
\[
 \log\frac{P_i^{\mathrm{old}}}{P_i^{\mathrm{new}}}
 =-\log F_i(M-\tau).
\]
It is nonnegative, and its conditional expected sum over all indices
outside the band is at most $13/12$ by \cref{eq:band-loss}.  The updated
mass of an actually hidden pair is positive almost surely; zero-mass
null events do not affect this expectation.

For any coordinate that is eventually force-exposed, its expenditures
up to that time telescope from initial mass one to a mass below
$3/4$.  Their sum therefore exceeds
\[
 \eta=\log(4/3).
\]
Pathwise, the number of forced exposures is at most the total
expenditure over all queries and coordinates divided by $\eta$.
Expenditure on coordinates that are later band-exposed only increases
this upper bound.  Summing the conditional expectation bounds and
adding the band exposures yields
\[
 \E R_T
 \le\frac{13}{12}T+\frac{13}{12\log(4/3)}T
 <5T.
\]
The argument allows arbitrary adaptive queries and includes the
independent algorithmic seed in the filtration.
\end{proof}

\begin{proposition}[A hidden sign prevents a uniformly correct output]
\label{prop:hidden-sign}
For $m\ge16$ and $T\le m/100$, every algorithm interacting with the
enhanced oracle satisfies
\[
 \Prob\bigl(\exists i\le m:\sigma_i\widehat x_i\ge0\bigr)
 \ge\frac34.
\]
\end{proposition}

\begin{proof}
By \cref{lem:budget},
$\E R_T<5T\le m/20$.  Markov's inequality gives
\[
 \Prob(R_T>m/2)\le\frac1{10}.
\]
Condition on the final enhanced transcript and the complete seed.  On
$R_T\le m/2$, at least $m/2$ pairs remain unexposed.  Their signs are
conditionally independent, and each sign value has probability at most
$2/3$ by \cref{lem:posterior}.  The output is fixed under this
conditioning.  If $\widehat x_i\ne0$, the inequality
$\sigma_i\widehat x_i<0$ requires one specified sign, whose conditional
probability is at most $2/3$.  If $\widehat x_i=0$, the inequality is
impossible.  Thus
\[
 \Prob(\sigma_i\widehat x_i<0\text{ for every }i\le m)
 \le\frac1{10}+\left(\frac23\right)^{m/2}
 \le\frac1{10}+\left(\frac23\right)^8
 <\frac14.
\]
Taking complements proves the claim.  The proof uses the actual output
directly, so it neither requires the output to have been queried nor
adds an extra oracle call.
\end{proof}

\subsection{The distributional lower bound}

\begin{proof}[Proof of \cref{thm:distribution}]
Sample the latent pairs from \cref{eq:latent} and define $h$ and $f$ by
\cref{eq:hard-max,eq:hard-smooth}.  This gives one distribution over
functions before the algorithm and its seed are considered.  By
\cref{eq:hard-radii}, \cref{lem:smoothing}, and
\cref{lem:norm-conversion}, every sampled function belongs to
$\cF_p(d,L,R_2,R_p)$.  Its global minimum is zero, attained at the point
in \cref{eq:hard-minimizer}, which satisfies both radius bounds.

Fix any randomized algorithm using at most $m/100$ exact first-order
calls.  By \cref{lem:simulation}, it can be run against the enhanced
oracle while ignoring the additional information; its simulated
transcript and output are exactly those obtained from the sampled
smooth function.  By \cref{prop:hidden-sign}, with probability at least
$3/4$ there is an index $i$ with $\sigma_i\widehat x_i\ge0$.  On this
event,
\[
 h(\widehat x)\ge\sigma_i\widehat x_i+b_i\ge a/2.
\]
The smoothing estimate in \cref{eq:smoothing-approx} then gives
\[
 (\cS_\mu h)(\widehat x)
 \ge h(\widehat x)-\mu/2
 \ge a/2-a/(256m)
 \ge a/3.
\]
Using \cref{eq:a-mu-tau,eq:hard-smooth}, we obtain
\begin{align*}
 f(\widehat x)-f^\star
 &\ge\frac{L\mu}{\kappa}\frac a3
 =\frac{La^2}{384\kappa m}\\
 &=\frac{L}{384\kappa}
   \min\left\{
       \frac{R_2^2}{m^2},
       \frac{R_p^2}{m^{1+2/p}}
   \right\}.
\end{align*}
This proves \cref{eq:distribution-lower}, with probability taken over
the sampled function and the independent seed.  The distribution was
fixed before choosing the algorithm, so the conclusion holds for every
algorithm with the stated budget.

Finally, averaging over this distribution shows that, for each fixed
randomized algorithm, some fixed instance has probability at least $3/4$ of an error at least
the displayed threshold.  Its expected error is at least $3/4$ times
that threshold.  The instance selected by averaging
may depend on the algorithm; the distribution does not.
\end{proof}

\section{From the hard distribution to the complexity frontier}
\label{app:main-theorem}

We give the parameter choices and integer arguments that connect the
fixed-distribution lower bound to the main complexity and risk statements.

\begin{lemma}[Uniform control of the smoothing factor]
\label{lem:kappa-uniform}
For $m\ge16$, take
\[
 s=\min\{p,\max\{2,\log m\}\}.
\]
Then the quantity in \cref{eq:kappa-def} satisfies
\[
 \kappa=(s-1)m^{2/s-2/p}\le e^2\log(em).
\]
\end{lemma}

\begin{proof}
Since $m\ge16$, we have $\log m>2$.
If $s=p$, then $\kappa=p-1\le\log(em)$.
Otherwise $s=\log m<p$, so
\[
 m^{2/s-2/p}\le m^{2/s}=e^2,
 \qquad s-1\le\log m.
\]
The same calculation includes $p=\infty$, with $1/p=0$.
\end{proof}

\begin{proof}[Proof of \cref{thm:main}]
The zero-query case is exactly \cref{prop:zero-query}.
Hence consider $0<\eps<Lr^2/2$.
The upper bound follows by selecting the least expensive of
\cref{prop:euclidean,prop:power,prop:dimension}, using the auxiliary
exponent \cref{eq:q-cap} for the uniform power-prox bound.
All three methods are deterministic, so their guarantees apply to
every criterion in \cref{def:criteria} and require no confidence
amplification.

For the lower bound, set $\alpha_p=1+2/p\in[1,2]$.
The error threshold in \cref{eq:distribution-lower} equals
\begin{equation}
\label{eq:threshold-rewrite}
 \frac{\eps}{384\kappa}
 \min\left\{
  \left(\frac A m\right)^2,
  \left(\frac B m\right)^{\alpha_p}
 \right\}.
\end{equation}
Choose the universal constants
\[
 K=3072e^2,
 \qquad c=\frac1{200K}.
\]
Suppose first that $S/(KH)\ge32$, and take
\[
 m=\left\lfloor\frac{S}{KH}\right\rfloor.
\]
Then $16\le m\le d$ and
\[
 m\ge\frac{S}{2KH},\qquad
 \kappa\le e^2H,\qquad
 \frac A m\ge KH,\qquad\frac B m\ge KH.
\]
Because $KH\ge1$ and $1\le\alpha_p\le2$, the threshold in
\cref{eq:threshold-rewrite} is at least
\[
 \frac{\eps}{384e^2H}(KH)=8\eps.
\]
By \cref{thm:distribution}, every algorithm using at most
$\lfloor m/100\rfloor$ calls has error at least $8\eps$ with probability
at least $3/4$ under the fixed hard distribution.
Its average expected error is therefore at least $6\eps$, and its
average failure probability at tolerance $\eps$ is at least $3/4$.
In each case, some fixed member of the distribution violates the
required guarantee.  This excludes deterministic algorithms as well,
since they are a special case of randomized algorithms.

The query counts are integers.  Thus, for every complexity $Q$ in the
statement,
\[
 Q\ge\lfloor m/100\rfloor+1
   >\frac m{100}
   \ge\frac{S}{200KH}
   =\frac{cS}{H}.
\]
This argument applies to every $0<\delta\le1/3$.

If instead $S/(KH)<32$, \cref{prop:zero-query} gives $Q\ge1$ in the
present accuracy range.  Since
\[
 \frac{cS}{H}<32cK=\frac{32}{200}<1,
\]
the same lower bound follows.
Together, the two cases prove \cref{eq:main-explicit}; its expression in
soft notation is \cref{eq:main-soft}.
\end{proof}

\begin{proof}[Proof of \cref{cor:risk}]
Fix finite $p\ge2$ and an integer $1\le T\le d/100$.
For the lower bound, choose $m=100T$ and $s=p$ in
\cref{thm:distribution}; then $m\le d$ and $\kappa=p-1$.
The distributional expected error of any randomized $T$-query
algorithm is at least
\[
 \frac{L}{512(p-1)}
 \min\left\{
  \frac{R_2^2}{(100T)^2},
  \frac{R_p^2}{(100T)^{1+2/p}}
 \right\}.
\]
A supremum over fixed instances is at least this average.
Since $1+2/p\le2$, the display is bounded below by
\[
 \frac{L}{512\cdot10^4(p-1)}
 \min\left\{
  \frac{R_2^2}{(T+1)^2},
  \frac{R_p^2}{(T+1)^{1+2/p}}
 \right\}.
\]
This gives an explicit admissible value of the constant $c_p$.

For the upper bound, \cref{prop:euclidean} gives
$2LR_2^2/(T+1)^2$.
Taking $q=p$ in \cref{prop:power} and choosing the target accuracy
\[
 \eps_T=LR_p^2\left(\frac{4p^2}{T}\right)^{1+2/p}
\]
gives an algorithm making $T$ calls with error at most $\eps_T$.
Because $(T+1)/T\le2$, this is at most
\[
 4(4p^2)^{1+2/p}\frac{LR_p^2}{(T+1)^{1+2/p}}.
\]
Selecting between these two algorithms using the known parameters
proves the deterministic upper bound with a constant $C_p$ depending
only on $p$.  Finally, deterministic algorithms are included among
randomized algorithms, so
$\Risk_T^{\mathrm{rand}}\le\Risk_T^{\mathrm{det}}$.
\end{proof}

\end{document}
